% ===================== page 1 (p001) =====================
\printedpage{}
% no page number printed (title page)
% head: (none)

\begin{center}
{\LARGE \textsf{THE DERIVATORS}}

\bigskip

{\large \textbf{ALEXANDRE GROTHENDIECK}}

\vspace{3em}

{\large \textbf{Chapter X}}

\medskip

{\large \textbf{Comparison of} Cat \textbf{with} $\bDelta^{\wedge}$}
\end{center}

\clearpage

% ===================== page 2 (p002) =====================
\printedpage{}
% no page number printed (editorial note)
% head: (none)

This text was deciphered and transcribed in \LaTeXe{} by M. Künzer. It was edited by M. Künzer, J. Malgoire, and G. Maltsiniotis. The transcription is as faithful as possible to the manuscript. For the few obvious corrections, or rare comments by the editors, as well as for the original numbering of the manuscript pages, typewriter characters \texttt{[typewriter]} in square brackets are used. A question mark in square brackets means that one is not sure of the preceding word.

This edition is provisional. Remarks, comments and corrections are welcome. Send a message to:

\textbf{[e-mail address removed]}

\hfill G. Maltsiniotis\hspace*{3em}

\clearpage

% ===================== page 3 (p003), printed 1 =====================
\printedpage{1}
% head: A. Grothendieck : Dérivateurs, Ch. X, Édition des Universités de Montpellier II et Paris VII
% foot: Version 9 avril 2002

\msp{page 1}

\section*{1\quad Recall of the notions of D-asphericity, and link with the ``basic localizer $W_{\mathbf{D}}$''}

For a derivator $\D$, the notions in question are the following (cf. I):

\begin{itemize}
\item[(i)] \textit{Object} of Cat that is $\D$-\textit{aspheric} (i.e. $Z \longrightarrow e$ is a $\D$-equivalence). NB I prefer to say ``$\D$-aspheric'' instead of ``$\D$-contractible'', which sometimes leads to confusion.

\item[\textcircled{ii}] $\D$-equivalence for an arrow $f : X \longrightarrow Y$ in Cat.

\item[(iii)] \textit{Functor} $f : X \longrightarrow Y$ that is $\D$-\textit{aspheric} (i.e. the $X/y$ are $\D$-aspheric) and $\D$-\textit{coaspheric} (the $y\backslash X$ are $\D$-aspheric). This is what I had called ``strict $\D$-equivalence'' (or ``costrict''), or homologically $\D$-cofinal functor (or cohomologically $\D$-cofinal), but the terminology ``aspheric'' seems to me the simplest and the best. (Fished out of PS, Pursuing Stacks.)

\item[(iv)] \textit{Functor} $f : X \longrightarrow Y$ that is $\D$-\textit{smooth,} $\D$-\textit{proper} (i.e. such that the fibers of the induced functors $X/x \longrightarrow Y/f(x)$, respectively $x\backslash X \longrightarrow f(x)\backslash Y$, are $\D$-aspheric).

\item[(v)] $\D$-\textit{fibration} $f : X \longrightarrow Y$ (self-dual notion), i.e. such that for $Y'' \xrightarrow{u} Y'$ over $Y$, if $Y''$, $Y'$ \texttt{[are]} Hot-aspheric, then $u_X : X'' \longrightarrow X'$ in
\[
\begin{tikzcd}
X \arrow[d] & X' \arrow[l] \arrow[d] & X'' \arrow[l] \arrow[d] \\
Y & Y' \arrow[l] & Y'' \arrow[l]
\end{tikzcd}
\]
is a $\D$-equivalence (\footnote{We shall also say that $f$ is $\D$-\textit{fibrant.} (\textbf{NB} There is no notion of ``$\D$-cofibrant'' to my knowledge.)}). The $\D$-\textit{fibration} is called \textit{trivial} (or \textit{aspheric}) if and only if its fibers are aspheric, or again if and only if $f$ is aspheric, or if and only if $f$ is coasperic. This also means that $f$ is a $\D$-equivalence that is \textit{universal,} % [sic] coaspérique -> coasphérique, letter h missing
\msp{page 2}
i.e. that for every $Y' \longrightarrow Y$, $X' = X \times_Y Y' \longrightarrow Y'$ is a $\D$-equivalence.

\item[(vi)] We also have weakened variants of the notion of $\D$-fibrations, namely the ``$\D$-\textit{left fibrations'',} or the $f$ that are ``\textit{cohomologically} $\D$-\textit{fibrant'',} and likewise the ``$\D$-\textit{right fibrations'',} i.e. the $f$ that are ``\textit{homologically} $\D$-\textit{fibrant''} (\footnote{New terminology: cohomological (resp. homological) fibrations that are \textit{weak,} cf. I, pages 111-126 for a corrected revision of the various notions of ``$\D$-fibrations''.}). The ``left''
\printedpage{2}
% head: A. Grothendieck : Dérivateurs, Ch. X, Édition des Universités de Montpellier II et Paris VII
% foot: Version 9 avril 2002
or cohomological notion for $f : X \longrightarrow Y$ means that the functors $X_y \hookrightarrow X/y$ ($y \in \Ob Y$) and $X/y \longrightarrow X/y'$ (for $(u : y \longrightarrow y') \in \Fl Y$) are $\D$-equivalences. And dually with the $X_y \hookrightarrow y\backslash X$, $y'\backslash X \longrightarrow y\backslash X$ for the ``right'' or ``homological'' notion. The $\D$-fibrations are none other than the universal left (cohomological) $\D$-fibrations or the universal right (homological) $\D$-fibrations.
\end{itemize}

Finally, this whole string of notions is deduced from the single notion of $\D$-equivalence (ii). If I denote by
\[
W_{\D} \subseteq \Fl(\Cat)
\]
the set of these $\D$-equivalences, one finds what in Pursuing Stacks I called a ``localizer'' (localizer) in Cat (\footnote{But beware, I assume here that the domain Diag of $\D$ is the whole of Cat!}). In Pursuing Stacks it was such a subset
\[
W \subseteq \Fl(\Cat),
\]
and not the datum of a derivator, that
\msp{page 3}
was taken as the fundamental starting datum for doing homotopical algebra of all kinds. I had moreover developed a string of notions of the same kind (i) to (vi), starting from this datum $W$. It could be that the statements for the notions I developed in I, in terms of a derivator, are valid more generally in terms of a localizer $W$, at least insofar as they still make sense. (So for the moment I set aside the strictly cohomological or homological statements, which appeal directly to $\D$.) Although it could be that under conditions to be made explicit on $W$, which imply that $(\Cat, W)$ do indeed define a derivator $\D_{\Cat,W} = \D$, one still manages to interpret $W$ as a $W_{\D}$, so that the cohomological and homological statements still make sense.

I feel like first looking in Pursuing Stacks at the axioms I have singled out for localizers, and checking whether they are all satisfied for
\msp{page 4}
$W_{\D}$. These conditions L 1 to L 5 are stated in Pursuing Stacks, §64 (page 164).

\textbf{L 1}\enspace Saturation of $W_D$. O.K.

This follows formally from the fact that we have a functor
\[
\Cat \xrightarrow{\varphi} C \qquad (= \Cat^o)
\]
such that for $u \in \Fl(\Cat)$, $u \in W_{\D} \iff \varphi(u)$ isomorphism. (Take for $\varphi(X)$ the subcategory of $\D(X)$ consisting of the $\D$-coefficients \textit{constant} on $X$).

\textbf{L 2}\enspace ``$\bDelta^1$ is $W$-aspheric over $e$'', i.e. for every $X$ in Cat, $X \times \bDelta^1 \longrightarrow X$ is in $W$.

This also means that $\bDelta^1 \longrightarrow e$ is a trivial $\D$-fibration, which is indeed the case.

\clearpage

% ===================== page 5 (p005), printed 3 =====================
\printedpage{3}
% head: A. Grothendieck : Dérivateurs, Ch. X, Édition des Universités de Montpellier II et Paris VII
% foot: Version 9 avril 2002

\textbf{L 3}\enspace If $X \in \Ob\Cat$ has a final object, then $X \longrightarrow e$ is in $W$.

In fact, we require for derivators that this be true assuming only $X$ aspheric (in the absolute sense), and one even concludes from this that $X \longrightarrow e$ is $W$-aspheric, i.e. $X \times Y \xrightarrow{\mathrm{pr}_2} Y$ is in $W$ for every $Y$ in Cat (\footnote{In fact L 2 $\Longrightarrow$ L 3.}).

\textbf{L 4}\enspace If $f : X \longrightarrow Y$ is such that $X/y \longrightarrow Y/y$ are in $W$, then $f$ is in $W$.

This is OK.

\msp{page 5}
\textbf{L 5}\enspace If $X$, $Y$ are Cat-fibered over $S$, and $X \longrightarrow Y$ is an $S$-morphism that is \textit{cartesian,} inducing on the fibers $X_s \longrightarrow Y_s$ which are in $W$, then $f$ is in $W$.

For $W_{\D}$, this is proposition 6, page 71 of I.

These axioms L 1 to L 5 (equivalent to L 1, L 3$'$ ($\bDelta^1 \longrightarrow e$ is in $W$) and L 5) are in the nature of stability axioms, ensuring that $W$ is large enough. There is another condition which limits $W$ to not being too large:

\textbf{L a)}\enspace If $f \in W$, then $\pi_0(f)$ is bijective.

This is not satisfied for all the $W_{\D}$ – e.g. if the derivator $\D$ is \textit{zero,} then $W_{\D} = \Fl\Cat$! But for the derivators useful in practice, this always seems to be satisfied.

\textbf{L b), c)}\enspace The functor
\[
\Cat \longrightarrow \Cat\, W^{-1}
\]
commutes with finite sums and finite products (\footnote{This works if $(\Cat, W)$ comes from a Quillen triple (with Q-fibrations, Q-cofibrations, Q-weak equivalences).}).

To examine what this means in concrete terms on $W$ – I have not checked that it is O.K. for the $W_{\D}$.

\textbf{L d)}\enspace If $f : X \longrightarrow Y$ in $W$, then $f^o : X^o \longrightarrow Y^o$ as well.

This is corollary 2, page 75 of I.